\textbf{Proof.}
At \(\xi_0\), thm:triangle-completion-separation-certificate gives the ordinary gap \(\delta>0\), optimizer-wide slack \(\eta>0\), and the strict singular-stratum gaps
\[
\frac{31}{80}-\frac{87}{400}=\frac{17}{100}>0,
\qquad
\frac{39}{80}-\frac{2613}{8432}=\frac{468}{2635}>0.                 \tag{1}
\]
We first show that the rank statements in (1) persist.  Total-variation convergence of the assignment law on its finite support makes \(\Pi\) and \(\Omega\) continuous.  Exposure positivity makes \(\Pi^{-1}\) continuous on a neighborhood of \(\xi_0\).  Wasserstein convergence of the bounded marginals gives convergence of their first and second moments, while lem:restricted-marginal-extrema-continuity gives convergence of the two covariance endpoints.  Def:moment-matrix then gives convergence of both endpoint moment matrices.  The triangle matrix \(\Omega(\xi_0)\) is positive definite, so, after restricting to a relative-open neighborhood, \(\Omega(\xi)\) remains uniformly positive definite.

Under the normalization in def:novel-pseudoinverse-completion, the rank-zero component is a singleton, each rank-one component is a radius-\(\sqrt n\) sphere modulo sign, and the rank-two component is a product of two such quotients.  These are compact components distinguished by which arm columns vanish.  When \(\Omega(\xi)\succ0\), the rank of \(S_B\) is constant on each component.  On rank zero the criteria are explicit constants; on rank one the Moore--Penrose formula is a scalar inverse with denominator bounded away from zero by uniform positive definiteness; and on rank two it is an ordinary inverse whose smallest eigenvalue is uniformly positive on the compact normalized component.  Hence def:completion-pseudoinverse-loss makes \(g_B^+(0)\) and \(w^+(B)\) jointly continuous on every component, uniformly on compact neighborhoods.  Taking maxima or minima over a fixed compact component preserves continuity of the componentwise optimum value.  The two strict inequalities in (1) therefore persist on a smaller relative-open neighborhood.  Every completion optimizer of either criterion consequently has rank two there.

Next we show that every such optimizer lies strictly inside the \(\kappa\)-regular ordinary class.  If this failed arbitrarily near \(\xi_0\), there would be \(\xi_j\to\xi_0\) and rank-two completion optimizers \(B_j\) for one of the two criteria such that
\[
\lambda_{\min}\{n^{-1}S_{B_j}(\xi_j)\}<\kappa+\frac\eta2.
\]
Compactness of the normalized rank-two component supplies a convergent subsequence \(B_j\to\bar B\).  Joint continuity, together with comparison to every fixed completion basis, shows that \(\bar B\) is an optimizer of the same criterion at \(\xi_0\).  Thm:triangle-completion-separation-certificate identifies every such optimizer with an ordinary optimizer and gives normalized eigenvalue at least \(\kappa+\eta\), contradicting the preceding display after passage to the limit.  Thus every nearby completion optimizer has normalized eigenvalue at least \(\kappa+\eta/2\) and has an ordinary preimage in \(\mathcal B_\kappa(\xi)\).

Lem:zhao-rank-two-stratum-identification preserves the two criteria under this preimage.  Because the global completion optimum is attained inside \(\mathcal B_\kappa(\xi)\), its value equals the corresponding optimum over \(\mathcal B_\kappa(\xi)\).  Both inclusions of optimizer sets follow, giving
\[
\mathcal A_0^+(\xi)=\iota_\xi\{\mathcal A_0(\xi)\},
\qquad
\mathcal A_{\mathrm{mm}}^+(\xi)=\iota_\xi\{\mathcal A_{\mathrm{mm}}(\xi)\}. \tag{2}
\]

Finally, the freshly established lem:dominance-gap-lower-semicontinuity and the exact gap at \(\xi_0\) imply, by the same sequential contradiction argument, that a still smaller relative-open neighborhood satisfies
\[
\Delta(\xi)>\frac\delta2=\frac{2456}{39525}.
\]
For every ordinary zero-covariance optimizer and every ordinary minimax optimizer, their loss difference is at least \(\Delta(\xi)\), hence exceeds \(\delta/2\).  Equation (2) and preservation of loss under \(\iota_\xi\) give the identical conclusion for every pair of completion optimizers. \(\square\)